\chapter{Problem-Solving Patterns}\label{ch:patterns}

Algorithms are tools; patterns are the habits of thought that tell you which tool to pick up. Experienced competitors do not
solve a new problem from scratch: they recognise that it ``is a contribution problem'', ``wants to be run backwards'', or
``has a small answer, so enumerate it''. None of these patterns is an algorithm you could put in a library, yet together
they solve a large fraction of all contest problems, from Bronze to the IOI.

Each section below describes one pattern, shows it on a concrete problem with a figure, and ends with the signals that
suggest it. Where a problem needs code, the code is complete and was tested against a brute force.

\begin{chapterintro}{\silver\ \gold\ \plat}{Chapter~\ref{ch:approach}; some examples use sorting, prefix sums and a Fenwick
tree, which later chapters explain in full.}
  \item count by \emph{contribution}, and count the complement;
  \item fix one part of the answer and solve the rest fast;
  \item reverse time, go offline, and sweep over events;
  \item prove impossibility with invariants and termination with monovariants;
  \item use pigeonhole, extremal and reflection arguments;
  \item exploit answers that are small, split heavy from light, and work with differences instead of values.
\end{chapterintro}

% ======================================================================================================================
\section[Contribution: count each piece once]{Contribution: count each piece once \hfill \gold}\label{sec:contribution}

Many counting problems ask for a sum over a huge family of objects: over all subarrays, all pairs, all paths in a tree.
Enumerating the objects is hopeless (there are $\Theta(n^2)$ subarrays), but the sum can be regrouped: instead of asking
``what does each object contribute?'', ask ``for each \emph{element}, in how many objects does it play a role, and what does
it contribute there?''. The second question usually has a closed-form answer. This regrouping is the \defn{contribution
technique}, and it is the single most useful idea in counting problems.

\begin{problem}{Sum of subarray minimums}
Given $a_0, \dots, a_{n-1}$ ($n \le 2\cdot10^5$, values up to $10^6$), compute the sum of $\min(a_l, \dots, a_r)$ over all
$\frac{n(n+1)}{2}$ subarrays.
\end{problem}

Fix an element $a_i$ and ask in which subarrays it is the minimum. Going left from $i$, the subarray may extend until it
would include an element smaller than $a_i$; going right, likewise. If $L_i$ is the index of the previous smaller element and
$R_i$ the index of the next smaller one, then $a_i$ is the minimum of exactly the subarrays that start in $(L_i, i]$ and end in
$[i, R_i)$, and there are $(i - L_i)(R_i - i)$ of them (Figure~\ref{fig:contribution}). So
\[
  \sum_{l \le r} \min(a_l, \dots, a_r) \;=\; \sum_{i} a_i \cdot (i - L_i)\cdot(R_i - i).
\]
Both $L_i$ and $R_i$ come from a monotonic stack in $\Oh{n}$ (Chapter~\ref{ch:prefix}). One detail needs care: when equal
values occur, a subarray with two equal minima must be counted once, not twice. Breaking ties asymmetrically does it: use
``strictly smaller'' on the left and ``smaller or equal'' on the right, so the subarray is charged to the last copy of its
minimum.

\begin{figure}[htbp]
\centering
\begin{tikzpicture}
  % bars of height 0.34 cm per unit; index i occupies x in [i, i + 0.8]
  \fill[fgsky!15, rounded corners=2pt] (1.9, -0.1) rectangle (4.9, 2.95);
  \foreach \v [count=\i from 0] in {5,2,6,4,7,3,8} {
    \edef\st{\ifnum\i=3 hot\else\ifnum\i=1 good\else\ifnum\i=5 good\else n\fi\fi\fi}
    \draw[fill=fgfill, draw=fgink!70, line width=0.6pt, \st] (\i, 0) rectangle ++(0.8, {\v*0.34});
    \node[font=\sffamily\small, above] at ({\i+0.4}, {\v*0.34}) {\v};
    \node[lbl, below] at ({\i+0.4}, 0) {\i};
  }
  % below the axis: the nearest smaller elements, then the choices of start and end
  \node[font=\sffamily\scriptsize\bfseries, text=fggreen!80!black] at (1.4, -0.85) {$L_3$};
  \node[font=\sffamily\scriptsize\bfseries, text=fggreen!80!black] at (5.4, -0.85) {$R_3$};
  \draw[mbrace] (2.0, -0.62) -- (3.8, -0.62);
  \node[font=\sffamily\scriptsize, text=fgverm] at (2.9, -1.3) {start $l \in \{2, 3\}$};
  \draw[mbrace] (3.0, -1.55) -- (4.8, -1.55);
  \node[font=\sffamily\scriptsize, text=fgverm] at (3.9, -2.23) {end $r \in \{3, 4\}$};
  \node[notel, anchor=west, text width=5.6cm] at (7.4, 1.2) {The element $4$ at index 3 is the minimum of the four
    subarrays $[6,4]$, $[6,4,7]$, $[4]$ and $[4,7]$: every choice of a start in $\{2,3\}$ and an end in $\{3,4\}$. Its
    window stops just before the nearest smaller elements $L_3 = 1$ and $R_3 = 5$. It contributes
    $4 \cdot 2 \cdot 2 = 16$.};
\end{tikzpicture}
\caption{The contribution of one element. It rules the shaded window, which stops just before the nearest smaller
elements on each side (green). Summing $a_i \cdot (\text{choices of start}) \cdot (\text{choices of end})$ over all $i$
counts every subarray exactly once.}
\label{fig:contribution}
\end{figure}

\begin{code}
void solve() {
    int n;
    cin >> n;
    vi a(n);
    cin >> a;
    // a[i] is the minimum of a[l..r] iff left[i] < l <= i <= r < right[i]
    vi left(n), right(n), st;
    rep(i, n) {  // previous strictly smaller element
        while (!st.empty() && a[st.back()] >= a[i]) st.pop_back();
        left[i] = st.empty() ? -1 : st.back();
        st.push_back(i);
    }
    st.clear();
    per(i, n) {  // next smaller or equal element: a tie goes to its last copy
        while (!st.empty() && a[st.back()] > a[i]) st.pop_back();
        right[i] = st.empty() ? n : st.back();
        st.push_back(i);
    }
    ll ans = 0;
    rep(i, n) ans += a[i] * (i - left[i]) * (right[i] - i);
    print(ans);
}
\end{code}

The same regrouping works far beyond minimums. The sum of $|a_i - a_j|$ over all pairs becomes, after sorting, $\sum_i
a_i\,(2i - n + 1)$: the $i$-th smallest value is added $i$ times (against every smaller value) and subtracted $n - 1 - i$
times. The sum over all paths of a tree of their lengths becomes a sum over edges of (size of one side) $\times$ (size of the
other). The expected number of fixed points of a random permutation is $n \cdot \frac1n = 1$ by linearity, without ever
looking at a single permutation.

\begin{recognize}
``Sum over all subarrays / pairs / paths / subsets of $f$'' where enumerating the objects is too slow. Ask for each element
(or edge, or value) how many objects it influences and how.
\end{recognize}

% ======================================================================================================================
\section[Fix one thing]{Fix one thing, then solve the rest fast \hfill \silver}\label{sec:fixone}

When an answer involves several choices (a triple of indices, a subarray's two ends, a path's two endpoints), fixing one of
them often makes the rest easy. The art is choosing \emph{which} one to fix: it is usually the one in the middle, the one
that plays a special role (the maximum, the root, the first move), or the one that makes the remaining choices independent.

\begin{problem}{Increasing triplets}
Count the triples $i < j < k$ with $a_i < a_j < a_k$ ($n \le 2 \cdot 10^5$).
\end{problem}

Fixing $i$ leaves a two-dimensional problem. Fixing the \emph{middle} index $j$ splits it into two independent
one-dimensional counts: the number of smaller elements to its left, and the number of larger elements to its right. Every
pair of a left partner and a right partner is a valid triple, so $j$ contributes the product (Figure~\ref{fig:fixmiddle}).
Both counts come from a Fenwick tree over compressed values (Chapter~\ref{ch:segtree}): sweep left to right counting
smaller elements seen so far, then right to left counting larger ones.

\begin{figure}[htbp]
\centering
\begin{tikzpicture}
  \arrs{A}{0}{0}{2/good, 5/dim, 3/hot, 6/cool, 4/cool}
  \idx{A}{5}
  \spanabove{A-1}{A-2}{1 smaller}
  \spanabove{A-4}{A-5}{2 larger}
  \node[font=\sffamily\scriptsize\bfseries, text=fgverm] at ($(A-3.north)+(0,0.3)$) {$j$};
  \node[notel, anchor=west, text width=6.2cm] at (4.2, -0.1) {With $a_j = 3$ in the middle, the triples are $(2, 3, 6)$
    and $(2, 3, 4)$: one choice on the left times two on the right.};
\end{tikzpicture}
\caption{Fixing the middle element makes the two sides independent: $j$ contributes (smaller on the left) $\times$ (larger
on the right).}
\label{fig:fixmiddle}
\end{figure}

\begin{code}
void solve() {
    int n;
    cin >> n;
    vi a(n);
    cin >> a;
    vi vals = a;  // compress values to 0..m-1
    sort(all(vals)), vals.erase(unique(all(vals)), vals.end());
    for (auto& x : a) x = lower_bound(all(vals), x) - vals.begin();
    int m = sz(vals);
    vi bit(m + 1, 0);
    auto add = [&](int i) { for (i++; i <= m; i += i & -i) bit[i]++; };
    auto count_below = [&](int v) {  // how many stored values are < v
        ll s = 0;
        for (; v > 0; v -= v & -v) s += bit[v];
        return s;
    };
    vi smaller_left(n);
    rep(j, n) smaller_left[j] = count_below(a[j]), add(a[j]);
    fill(all(bit), 0);
    ll ans = 0;
    per(j, n) {  // fix the middle element j
        ll larger_right = (n - 1 - j) - count_below(a[j] + 1);
        ans += smaller_left[j] * larger_right;
        add(a[j]);
    }
    print(ans);
}
\end{code}

Other classic ``fix one thing'' moves: fix the right end of a subarray and maintain something about all left ends (prefix
sums with a hash map, Chapter~\ref{ch:prefix}); fix the maximum of a subarray and count the subarrays where it is the
maximum (the contribution technique again); fix the lowest common ancestor of a tree path (Part~IV); fix the value of the
answer and check it (binary search on the answer, Chapter~\ref{ch:binsearch}); fix the number of items taken from one group
and optimise the rest.

\begin{recognize}
Triples, pairs or intervals with a condition that couples all parts. Look for the part whose choice makes the others
independent or one-dimensional.
\end{recognize}

% ======================================================================================================================
\section[Complementary counting]{Complementary counting \hfill \silver}\label{sec:complement}

Counting objects that satisfy ``at least one'', ``not all'' or ``some pair differs'' is often awkward, while their complement
(``none'', ``all'', ``all equal'') is simple. Then count the complement and subtract it from the total. It sounds trivial,
and it is, but it is missed surprisingly often under contest pressure.

\begin{problem}{Subarrays with a large element}
Count the subarrays that contain at least one element $\ge K$.
\end{problem}

A subarray \emph{without} a large element lies entirely inside a maximal run of small elements, and a run of length $\ell$
contains $\frac{\ell(\ell+1)}{2}$ subarrays. So the answer is $\frac{n(n+1)}{2}$ minus the sum over runs
(Figure~\ref{fig:complement}), computed in one pass.

\begin{figure}[htbp]
\centering
\begin{tikzpicture}
  \arrs{A}{0}{0}{3/n, 7/hot, 2/n, 1/n, 8/hot, 4/n}
  \cellframe{A-1}{A-1}{fgblue} \cellframe{A-3}{A-4}{fgblue} \cellframe{A-6}{A-6}{fgblue}
  \foreach \c in {A-1, A-6} \node[font=\sffamily\scriptsize, text=fgblue, below=4pt] at (\c.south) {$\ell{=}1$: 1};
  \node[font=\sffamily\scriptsize, text=fgblue, below=4pt] at ($(A-3.south)!0.5!(A-4.south)$) {$\ell{=}2$: 3};
  \node[notel, anchor=west, text width=6.4cm] at (4.6, 0) {$K = 5$. All subarrays: $\frac{6 \cdot 7}{2} = 21$. Without a
    large element (inside a blue run): $1 + 3 + 1 = 5$. So $21 - 5 = 16$ subarrays contain a 7 or an 8.};
\end{tikzpicture}
\caption{Counting the complement: subarrays that avoid the large elements (orange) live inside the maximal runs of small
elements (blue), and a run of length $\ell$ holds $\ell(\ell+1)/2$ of them.}
\label{fig:complement}
\end{figure}

Complementary counting generalises to \defn{inclusion--exclusion} when there are several ``bad'' events (Chapter
\ref{ch:combinatorics}): the number of objects avoiding all bad events is the total, minus those with each single bad
event, plus those with each pair, and so on. It is also the natural way to count pairs of \emph{distinct} things: count
all pairs, subtract the pairs that are equal.

\begin{pitfall}[Double subtraction]
When the ``bad'' objects overlap (a subarray can contain two different forbidden patterns), subtracting each kind separately
removes some objects twice. Either make the bad cases disjoint by definition, as the maximal runs above are, or use
inclusion--exclusion.
\end{pitfall}

% ======================================================================================================================
\section[Reverse time and go offline]{Reverse time and go offline \hfill \gold}\label{sec:reverse}

Some operations are easy in one direction and hard in the other. A disjoint set union merges components in almost constant
time but cannot split them; a set supports insertion and deletion, but a max-heap cannot delete an arbitrary element. When
the problem gives you every operation in advance, you can process them in whatever order is convenient, including
\emph{backwards}. Deletions then become insertions.

\begin{problem}{Network Breakdown (CSES)}
A network has $n$ computers and $m$ connections. Connections are removed one at a time, $k$ times; after each removal, print
the number of connected components. Limits: $n \le 10^5$, $m, k \le 2 \cdot 10^5$.
\end{problem}

Forwards, every removal can split a component, and nothing simple maintains that. Backwards, start from the final network
(all $k$ connections already removed), compute its components with a DSU, then add the removed connections back in reverse
order. The number of components before adding connection $i$ back is exactly the answer after removal $i$
(Figure~\ref{fig:reverse}).

\begin{figure}[htbp]
\centering
\begin{tikzpicture}[mini/.style={vtxs, minimum size=4.5mm}]
  \newcommand{\ministate}[4]{% x offset, edges to draw, label below, components
    \begin{scope}[xshift=#1 cm]
      \node[mini] (a) at (0, 1.0) {1}; \node[mini] (b) at (1.0, 1.0) {2};
      \node[mini] (c) at (1.0, 0) {3}; \node[mini] (d) at (0, 0) {4};
      \foreach \u/\v in {#2} \draw[ed] (\u) -- (\v);
      \node[lbl] at (0.5, -0.6) {#3};
      \node[font=\sffamily\small\bfseries, text=fgink] at (0.5, -1.05) {#4};
    \end{scope}}
  \ministate{0}{a/b, b/c, c/d, a/d, a/c}{start}{1 comp.}
  \ministate{3.2}{a/b, b/c, c/d, a/d}{after $-\{1,3\}$}{1}
  \ministate{6.4}{a/b, c/d, a/d}{after $-\{2,3\}$}{1}
  \ministate{9.6}{a/b, c/d}{after $-\{1,4\}$}{2}
  \foreach \x in {1.4, 4.6, 7.8} {
    \draw[->, line width=0.8pt, fggray, dashed] (\x, 0.5) -- ++(1.4, 0);
    \draw[gded] ({\x+1.4}, -1.5) -- ++(-1.4, 0);
  }
  \node[font=\sffamily\scriptsize, text=fggray] at (6.05, 1.65) {forwards: each step removes an edge (a DSU cannot split)};
  \node[font=\sffamily\scriptsize, text=fggreen!80!black] at (6.05, -1.95) {backwards: each step adds an edge (a DSU
    merges in $\approx\Oh{1}$)};
\end{tikzpicture}
\caption{Forwards, the network loses edges; backwards, it gains them. Starting from the last state and adding the edges
back in reverse order visits the same states, so the component counts can be read off the DSU. Output: \texttt{1 1 2}.}
\label{fig:reverse}
\end{figure}

\begin{code}
void solve() {
    int n, m, k;
    cin >> n >> m >> k;
    set<pii> edges;
    rep(i, m) {
        int a, b;
        cin >> a >> b;
        edges.insert({min(a, b), max(a, b)});
    }
    vector<pii> removed(k);
    for (auto& [a, b] : removed) {
        cin >> a >> b;
        if (a > b) swap(a, b);
        edges.erase({a, b});
    }
    vi par(n + 1);
    iota(all(par), 0);
    auto root = [&](int x) {
        while (x != par[x]) x = par[x] = par[par[x]];
        return x;
    };
    int comps = n;
    auto unite = [&](int a, int b) {
        a = root(a), b = root(b);
        if (a != b) par[a] = b, comps--;
    };
    for (auto [a, b] : edges) unite(a, b);  // the graph after every removal
    vi ans(k);
    per(i, k) {  // play the removals backwards: each one becomes an insertion
        ans[i] = comps;
        unite(removed[i].first, removed[i].second);
    }
    print(ans);
}
\end{code}

\defn{Offline processing} is the general idea: when all queries are known before you must answer any of them, you may
reorder them. Sort queries by their right endpoint and sweep the array once (distinct values in ranges,
Chapter~\ref{ch:advds}); sort them into blocks and move two pointers (Mo's algorithm); answer them in a divide-and-conquer
over time (CDQ); or process them in reverse. Online problems, where each query depends on the previous answer (often through
an XOR with the last answer), deliberately forbid this.

\begin{recognize}
Deletions that would be insertions if time ran backwards; queries that are all given before any answer is needed; an
operation your data structure supports in only one direction.
\end{recognize}

% ======================================================================================================================
\section[Events and sweeps]{Events and sweeps \hfill \silver}\label{sec:events}

Problems about intervals in time or space become simple when you stop thinking about the intervals and start thinking
about their \emph{endpoints}. Each interval becomes two events, ``something starts'' and ``something ends''; sort the events
and walk through them in order, maintaining the current state. This is a \defn{sweep}, and in two dimensions it becomes the
sweep line of computational geometry (Chapter~\ref{ch:geometry}).

\begin{problem}{Restaurant Customers (CSES)}
Customer $i$ is in the restaurant from time $a_i$ to $b_i$. What is the largest number of customers present at the same
time? Limits: $n \le 2\cdot10^5$, all times distinct.
\end{problem}

Turn each customer into an arrival ($+1$) and a departure ($-1$), sort the $2n$ events by time, and keep a running sum.
The running sum is the number of customers present, and its maximum is the answer (Figure~\ref{fig:sweep}). Sorting costs
$\Oh{n \log n}$; the walk is linear. If equal times were allowed, the order of simultaneous events would matter: to count
someone who leaves at time $t$ and someone who arrives at $t$ as \emph{not} overlapping, process departures first. Sorting
pairs $(t, \pm 1)$ does exactly that, since $-1 < +1$.

\begin{figure}[htbp]
\centering
\begin{tikzpicture}[x=0.8cm]
  \foreach \a/\b/\y in {2/5/3.3, 3/8/2.75, 4/6/2.2, 7/9/1.65} \ivl{\a}{\b}{\y}{fgblue!55}{};
  \draw[->, line width=0.6pt, fgink!70] (0, 0) -- (10.5, 0) node[right, lbl] {time};
  \draw[->, line width=0.6pt, fgink!70] (0, 0) -- (0, 1.25) node[above, lbl] {present};
  \foreach \t in {0,...,10} \draw[line width=0.5pt, fgink!70] (\t, 0.06) -- (\t, -0.06) node[below, lbl] {\t};
  \draw[line width=1.4pt, fgverm] (0,0) -- (2,0) -- (2,0.3) -- (3,0.3) -- (3,0.6) -- (4,0.6) -- (4,0.9) -- (5,0.9) --
    (5,0.6) -- (6,0.6) -- (6,0.3) -- (7,0.3) -- (7,0.6) -- (8,0.6) -- (8,0.3) -- (9,0.3) -- (9,0) -- (10,0);
  \foreach \y/\v in {0.3/1, 0.6/2, 0.9/3} \node[lbl, anchor=east] at (-0.1, \y) {\v};
  \node[font=\sffamily\scriptsize\bfseries, text=fgverm] at (4.5, 1.2) {max 3};
  \draw[xed] (4, 0.9) -- (4, 3.5); \draw[xed] (5, 0.9) -- (5, 3.5);
\end{tikzpicture}
\caption{Four customers (bars) and the number present over time (the step function below). Each arrival lifts the step
by one and each departure lowers it; the maximum, 3, is reached between times 4 and 5.}
\label{fig:sweep}
\end{figure}

\begin{code}
void solve() {
    int n;
    cin >> n;
    vector<pii> ev;  // (time, +1 for an arrival, -1 for a departure)
    rep(i, n) {
        int a, b;
        cin >> a >> b;
        ev.push_back({a, 1}), ev.push_back({b, -1});
    }
    sort(all(ev));
    int cur = 0, best = 0;
    for (auto [t, d] : ev) chmax(best, cur += d);
    print(best);
}
\end{code}

\begin{trick}[Coordinate compression plus a difference array]
When the times are small integers (say up to $10^6$), skip the sort: add $+1$ at $a_i$ and $-1$ at $b_i$ in an array and
take prefix sums (a difference array, Chapter~\ref{ch:prefix}). When they are large but few, compress them first. The sweep
and the difference array are the same idea seen from two sides.
\end{trick}

% ======================================================================================================================
\section[Invariants and colourings]{Invariants and colourings \hfill \silver}\label{sec:invariants}

An \defn{invariant} is a quantity that an operation never changes. If the start and the goal have different values of an
invariant, the goal is unreachable, no matter how clever the sequence of operations. Invariants are how contest problems
prove that something is \emph{impossible}, and many ``print \texttt{NO} or a construction'' problems are really ``find the
invariant, then show everything else is reachable''.

\begin{problem}{The mutilated chessboard}
Remove two opposite corner squares from an $8 \times 8$ board. Can the remaining 62 squares be tiled by 31 dominoes?
\solution Colour the board like a chessboard (Figure~\ref{fig:chessboard}). Every domino covers one dark and one light square,
wherever it is placed: that is the invariant. Opposite corners have the same colour, so the mutilated board has 30 squares of
one colour and 32 of the other. Thirty-one dominoes would need 31 of each. Impossible.
\end{problem}

\begin{figure}[htbp]
\centering
\begin{tikzpicture}[s/.style={minimum size=4.6mm, inner sep=0pt}]
  \foreach \r in {0,...,7} \foreach \c in {0,...,7} {
    \pgfmathtruncatemacro{\cpar}{mod(\r+\c,2)}
    \ifnum\cpar=0 \fill[fgink!55] (\c*0.46, -\r*0.46) rectangle ++(0.46, -0.46);
    \else \fill[fgfill] (\c*0.46, -\r*0.46) rectangle ++(0.46, -0.46); \fi
  }
  \draw[line width=0.6pt, fgink!70] (0, 0) rectangle (3.68, -3.68);
  \foreach \r/\c in {0/0, 7/7} {
    \fill[white] (\c*0.46, -\r*0.46) rectangle ++(0.46, -0.46);
    \draw[fgverm, line width=1.2pt] (\c*0.46+0.08, -\r*0.46-0.08) -- ++(0.3, -0.3) (\c*0.46+0.38, -\r*0.46-0.08) -- ++(-0.3, -0.3);
  }
  \draw[line width=2pt, fgorange, rounded corners=2pt] (2*0.46+0.04, -3*0.46-0.04) rectangle (4*0.46-0.04, -4*0.46+0.04);
  \draw[line width=2pt, fgorange, rounded corners=2pt] (5*0.46+0.04, -5*0.46-0.04) rectangle (6*0.46-0.04, -7*0.46+0.04);
  \node[notel, anchor=west, text width=6.4cm] at (4.3, -1.84) {Every domino (orange) covers one dark and one light square,
    horizontal or vertical. Both removed corners (red crosses) are dark, leaving 30 dark and 32 light squares: no tiling
    can exist.};
\end{tikzpicture}
\caption{A colouring invariant. The count ``dark squares covered minus light squares covered'' is zero for any set of
dominoes, but it would have to be $-2$ here.}
\label{fig:chessboard}
\end{figure}

Parity is the most common invariant. If an operation swaps two elements of a permutation, the parity of the number of
inversions flips every time; so whether a target permutation is reachable in an even number of swaps is decided before you
start (this is why half the configurations of the 15-puzzle are unsolvable). If an operation replaces two numbers by their
sum and difference, the parity of the total sum never changes. Other common invariants are a sum modulo some $m$, the
greatest common divisor of all elements, the number of connected components, and weighted sums like $\sum_i a_i \cdot 2^i$.

\begin{tip}[Finding the invariant]
Apply the operation to small examples and write down what stays the same: the sum, its parity, the sum mod 3, the
multiset of values, the gcd, the number of odd elements. When the problem colours something, try colouring it yourself:
chessboard colouring for dominoes, colouring by $x \bmod 3$ or $(x + y) \bmod k$ for longer tiles.
\end{tip}

% ======================================================================================================================
\section[Monovariants and termination]{Monovariants and termination \hfill \gold}\label{sec:monovariant}

A \defn{monovariant} is a quantity that changes in only one direction: every operation strictly decreases it, say, and it
cannot go below zero. Then the process must stop, and the number of steps is at most the starting value. Monovariants are
how you prove that a local-improvement algorithm terminates, and they bound its running time.

A classic example: ``while some adjacent pair is out of order, swap it''. Each swap removes exactly one inversion (a pair
$i < j$ with $a_i > a_j$), and the number of inversions is at most $\binom n2$, so the process ends after exactly
\emph{inversions} swaps. That count is also the minimum number of adjacent swaps that sorts the array, because no single
swap can remove more than one inversion. Counting inversions quickly (merge sort or a Fenwick tree, $\Oh{n\log n}$) thus
answers ``how many adjacent swaps are needed to sort?''.

Monovariants also justify local search. ``While some vertex has more than half of its neighbours on its own side, move it
to the other side'' terminates, because each move strictly increases the number of edges between the two sides, which is at
most $m$. When it stops, every vertex has at least half of its neighbours across, which proves that a cut of at least $m/2$
edges always exists, and finds one in at most $m$ moves.

\begin{recognize}
``Show that the process ends'', ``how many operations until nothing changes'', or a local-improvement algorithm whose
running time you need to bound. Look for something that strictly improves: a sum, a count of bad pairs, a potential.
\end{recognize}

% ======================================================================================================================
\section[Pigeonhole]{Pigeonhole \hfill \silver}\label{sec:pigeonhole}

The \defn{pigeonhole principle}: if $n + 1$ objects go into $n$ boxes, some box holds two of them. It is as obvious as it
sounds and as powerful as any theorem in this book, because it proves that something \emph{exists} without constructing
it, and the proof usually shows how to find it.

\begin{problem}{Subarray Divisibility (CSES)}
Count the subarrays of $a_1, \dots, a_n$ whose sum is divisible by $n$ ($n \le 2\cdot10^5$, values may be negative).
\solution The subarray $a_{l+1}, \dots, a_r$ has sum $P_r - P_l$, which is divisible by $n$ exactly when the prefix sums
$P_l$ and $P_r$ have the same remainder mod $n$. Put the $n + 1$ prefix sums $P_0, \dots, P_n$ into $n$ boxes by remainder
(Figure~\ref{fig:pigeonhole}); every pair inside one box is a good subarray. The pigeonhole principle adds a bonus: since
there are $n + 1$ prefix sums and only $n$ remainders, \emph{some} box has two, so a divisible subarray always exists.
\end{problem}

\begin{figure}[htbp]
\centering
\begin{tikzpicture}
  % a = 3 1 4 1 5, prefix sums P_0..P_5 = 0 3 4 8 9 14, remainders 0 3 4 3 4 4
  \arr{A}{0.35}{1.1}{3,1,4,1,5}
  \node[lbl, anchor=east] at (0.2, 1.1) {$a$};
  \arr{P}{0}{0}{0,3,4,8,9,14}
  \node[lbl, anchor=east] at (-0.15, 0) {$P_0..P_5$};
  \arrs{M}{0}{-0.85}{0/cool, 3/good, 4/hot, 3/good, 4/hot, 4/hot}
  \node[lbl, anchor=east] at (-0.15, -0.85) {$P \bmod 5$};
  \foreach \r/\x in {0/5.0, 1/6.3, 2/7.6, 3/8.9, 4/10.2} {
    \draw[line width=0.7pt, fgink!70] (\x, 1.1) -- (\x, -0.95) -- (\x+1.1, -0.95) -- (\x+1.1, 1.1);
    \node[lbl] at (\x+0.55, -1.2) {rem.\ \r};
  }
  \node[vtxs, cool, minimum size=5.5mm] at (5.55, -0.6) {$P_0$};
  \node[vtxs, good, minimum size=5.5mm] at (9.45, -0.6) {$P_1$};
  \node[vtxs, good, minimum size=5.5mm] at (9.45, 0.05) {$P_3$};
  \node[vtxs, hot, minimum size=5.5mm] at (10.75, -0.6) {$P_2$};
  \node[vtxs, hot, minimum size=5.5mm] at (10.75, 0.05) {$P_4$};
  \node[vtxs, hot, minimum size=5.5mm] at (10.75, 0.7) {$P_5$};
\end{tikzpicture}
\caption{$n = 5$: six prefix sums and five possible remainders, so some box must hold two. Box 3 holds $P_1, P_3$ (the
subarray $1, 4$); box 4 holds $P_2, P_4, P_5$ (the subarrays $4, 1$ and $4, 1, 5$ and $5$). In total
$\binom22 + \binom32 = 4$ subarrays have a sum divisible by 5.}
\label{fig:pigeonhole}
\end{figure}

\begin{code}
void solve() {
    int n;
    cin >> n;
    vi cnt(n, 0);  // how many prefixes so far have each remainder mod n
    cnt[0] = 1;
    ll pre = 0, ans = 0;
    rep(i, n) {
        ll x;
        cin >> x;
        pre = ((pre + x) % n + n) % n;
        ans += cnt[pre]++;  // pair with earlier prefixes of the same remainder
    }
    print(ans);
}
\end{code}

Pigeonhole arguments also bound search spaces. A sequence defined by $x_{i+1} = f(x_i)$ over $m$ possible values repeats
within $m + 1$ steps, so it is eventually periodic (Chapter~\ref{ch:graphs}, functional graphs). Among any $n + 1$ integers
in $[1, 2n]$, two are coprime (two are consecutive). Among 5 points in a $2 \times 2$ square, two are at distance at most
$\sqrt2$. Whenever a problem says ``prove that there always exists'', look for boxes and pigeons.

% ======================================================================================================================
\section[The extremal principle]{The extremal principle \hfill \silver}\label{sec:extremal}

\emph{Look at the largest, the smallest, the first or the last.} An extreme object has properties that ordinary objects do
not, because nothing can exceed it. Many existence proofs, and many greedy algorithms, start by picking an extreme object
and showing it must behave nicely.

\begin{problem}{Every tree has two leaves}
Show that every tree with at least two vertices has at least two leaves (vertices of degree 1).
\solution Take a \emph{longest} path in the tree, with endpoints $u$ and $v$ (Figure~\ref{fig:extremal}). If $u$ had a
neighbour $w$ other than the next vertex on the path, then either $w$ is on the path (and the tree would contain a cycle) or
it is not (and prepending $w$ would make a longer path). Both are impossible, so $u$ has degree 1; the same holds for $v$.
\end{problem}

\begin{figure}[htbp]
\centering
\begin{tikzpicture}
  \node[vtx, good] (u) at (0, 0) {$u$};
  \node[vtx] (a) at (1.4, 0) {}; \node[vtx] (b) at (2.8, 0) {}; \node[vtx] (c) at (4.2, 0) {};
  \node[vtx, good] (v) at (5.6, 0) {$v$};
  \node[vtx] (x) at (2.8, 1.2) {}; \node[vtx] (y) at (4.2, -1.2) {}; \node[vtx] (z) at (1.4, -1.2) {};
  \draw[hed] (u) -- (a) -- (b) -- (c) -- (v);
  \draw[ed] (b) -- (x) (c) -- (y) (a) -- (z);
  \node[vtx, dim, dashed] (w) at (-1.4, 0.9) {$w$};
  \draw[xed] (u) -- (w);
  \node[notel, anchor=west, text width=4.6cm] at (6.5, 0) {If $u$ had another neighbour $w$, the path $w, u, \dots, v$
    would be longer than the longest path. So $u$ and $v$ are leaves.};
\end{tikzpicture}
\caption{The extremal principle: the endpoints of a longest path cannot have any other neighbour.}
\label{fig:extremal}
\end{figure}

The same move appears everywhere in this book. The leftmost point of a point set is on its convex hull. The interval that
ends first is in some optimal schedule (Section~\ref{sec:exchange}). In a graph where every vertex has degree at least $d$,
a longest path's endpoint has all its neighbours on the path, so the graph has a cycle of length at least $d + 1$. The
smallest counterexample to a claim (``consider a minimal counterexample'') is the extremal principle applied to proofs by
contradiction, and it is exactly how induction proofs are usually discovered.

% ======================================================================================================================
\section[Symmetry and reflection]{Symmetry and reflection \hfill \gold}\label{sec:symmetry}

Symmetry lets you count one representative and multiply, or count everything and divide. Unordered pairs are ordered pairs
divided by two. ``Assume without loss of generality that $a \le b$'' halves the casework. Rotations of a necklace are
counted by Burnside's lemma (Chapter~\ref{ch:combinatorics}). The most beautiful use of symmetry in contest counting is the
\defn{reflection principle}.

\begin{problem}{Balanced bracket sequences}
How many sequences of $n$ opening and $n$ closing brackets are balanced (no prefix has more closing than opening brackets)?
\solution Draw a sequence as a lattice path from $(0, 0)$ to $(n, n)$: an opening bracket is a step right, a closing one a
step up. The sequence is balanced exactly when the path never goes above the diagonal, i.e.\ never touches the line
$y = x + 1$. Count the bad paths: reflect the part of a bad path \emph{after its first touch} of $y = x + 1$ across that
line (Figure~\ref{fig:reflection}). The result is a path to $(n - 1, n + 1)$, and every path to $(n - 1, n + 1)$ arises from
exactly one bad path this way (it must cross the line, and reflecting back undoes the map). So there are $\binom{2n}{n-1}$
bad paths, and the balanced ones number
\[ \binom{2n}{n} - \binom{2n}{n-1} = \frac{1}{n + 1}\binom{2n}{n}, \]
the Catalan number $C_n$.
\end{problem}

\begin{figure}[htbp]
\centering
\begin{tikzpicture}[scale=0.78]
  \draw[line width=0.3pt, fgline!60] (0, 0) grid (5, 5);
  \draw[line width=0.8pt, fgink!40] (0, 0) -- (5, 5);
  \draw[line width=1pt, fgverm, dashed] (0, 1) -- (4, 5);
  \node[font=\sffamily\scriptsize, text=fgverm, anchor=south east] at (0, 1.05) {$y = x + 1$};
  \draw[line width=2.2pt, fgblue] (0,0) -- (1,0) -- (1,1) -- (1,2) -- (2,2) -- (2,3) -- (3,3) -- (4,3) -- (4,4);
  \draw[line width=2.2pt, fgorange, dash pattern=on 4pt off 2pt] (1,2) -- (1,3) -- (2,3) -- (2,4) -- (2,5) -- (3,5);
  \node[circle, fill=white, draw=fgink, line width=1pt, minimum size=2.2mm, inner sep=0pt] at (1, 2) {};
  \draw[line width=0.5pt, fgink!60] (0.88, 2) -- (-0.35, 2.55);
  \node[font=\sffamily\scriptsize, anchor=east] at (-0.35, 2.55) {first touch};
  \node[circle, fill=fgblue, minimum size=1.8mm, inner sep=0pt] at (4, 4) {};
  \node[font=\sffamily\scriptsize, text=fgblue, anchor=west] at (4.15, 4) {$(4, 4)$};
  \node[circle, fill=fgorange, minimum size=1.8mm, inner sep=0pt] at (3, 5) {};
  \node[font=\sffamily\scriptsize, text=fgorange!80!black, anchor=west] at (3.15, 5.2) {$(3, 5)$};
  \node[lbl, anchor=north] at (2.5, -0.15) {right = \texttt{(}, up = \texttt{)}};
  \node[notel, anchor=west, text width=5.6cm] at (6.3, 2.5) {Blue: the bad path \texttt{( ) ) ( ) ( ( )}, which touches
    $y = x + 1$ (red, dashed) at the white dot. Swapping right and up steps after that point reflects the rest of the path
    (orange) and moves its end from $(4, 4)$ to $(3, 5)$. Every path to $(3, 5)$ arises this way exactly once.};
\end{tikzpicture}
\caption{The reflection principle. Reflecting the tail of a bad path across $y = x + 1$ moves its endpoint from $(n, n)$ to
$(n - 1, n + 1)$, which turns ``count the bad paths'' into a single binomial coefficient.}
\label{fig:reflection}
\end{figure}

\begin{fact}
The reflection principle counts lattice paths that avoid a line for any start and end point, and it solves ballot problems
(``candidate A is always strictly ahead during the count'') and bounded random walks. With several forbidden lines it
becomes an alternating sum of reflections.
\end{fact}

% ======================================================================================================================
\section[When the answer is small]{When the answer is small, enumerate it \hfill \plat}\label{sec:small-answer}

Sometimes a quantity that looks like it could be huge is provably tiny, and then you can simply list all its values. The
classic source is \emph{divisibility}: a chain of numbers where each divides the previous one can only have $\log_2 A + 1$
distinct values below $A$, because each strict step at least halves the value.

\begin{problem}{Subarrays with gcd 1}
Count the subarrays whose greatest common divisor is 1 ($n \le 2 \cdot 10^5$, values up to $10^9$).
\solution Fix the right end $i$ and look at the gcds of $a_l, \dots, a_i$ as $l$ moves left: each is a divisor of the
previous one, so there are at most $\log_2 10^9 + 1 \approx 31$ distinct values (Figure~\ref{fig:gcdchain}). Store them as
(value, count) pairs; moving to $i + 1$ means taking the gcd of every stored value with $a_{i+1}$ and merging equal
neighbours. The total cost is $\Oh{n \log A}$ gcd computations instead of $\Oh{n^2}$ subarrays.
\end{problem}

\begin{figure}[htbp]
\centering
\begin{tikzpicture}[g/.style={cellc, minimum width=8mm, minimum height=6.5mm, font=\sffamily\small}]
  \foreach \v [count=\c from 0] in {12, 18, 6, 9, 3} \node[hdr] at (\c*0.8, 0.75) {\v};
  \node[hdr, anchor=east] at (-0.55, 0.75) {$a$};
  \foreach \r/\row in {0/{12}, 1/{6,18}, 2/{6,6,6}, 3/{3,3,3,9}, 4/{3,3,3,3,3}} {
    \node[hdr, anchor=east] at (-0.55, -\r*0.65) {end $\r$};
    \foreach \v [count=\c from 0] in \row {
      \edef\st{\ifnum\v=3 good\else\ifnum\v=6 cool\else\ifnum\v=9 warn\else pur\fi\fi\fi}
      \node[g, \st] at (\c*0.8, -\r*0.65) {\v};
    }
  }
  \node[notel, anchor=west, text width=6.3cm] at (4.6, -1.3) {Row $i$ lists $\gcd(a_l, \dots, a_i)$ for every start $l$.
    Reading any row from right to left, each value divides the one before it, so a row has at most $\log_2 A + 1$
    \emph{distinct} values, however long it is. Row 3 is the chain $9 \to 3$ with multiplicities $1, 3$.};
\end{tikzpicture}
\caption{Subarray gcds of $12, 18, 6, 9, 3$, one row per right end. Colours mark equal values; every row is a short chain of
divisors, which is what makes the $\Oh{n \log A}$ algorithm possible.}
\label{fig:gcdchain}
\end{figure}

\begin{code}
void solve() {
    int n;
    cin >> n;
    vi a(n);
    cin >> a;
    vector<pll> cur;  // (gcd, count) for subarrays ending at i: a short chain
    ll ones = 0;
    rep(i, n) {
        vector<pll> nxt{{a[i], 1}};
        for (auto [g, c] : cur) {
            ll h = gcd(g, a[i]);
            if (nxt.back().first == h) nxt.back().second += c;
            else nxt.push_back({h, c});
        }
        cur = nxt;
        for (auto [g, c] : cur)
            if (g == 1) ones += c;
    }
    print(ones);
}
\end{code}

The same argument works for bitwise OR and AND over subarrays (each new element can only add or remove bits, at most 30
times), for the number of distinct values of $\lfloor n / i \rfloor$ (at most $2\sqrt n$, Chapter~\ref{ch:numbertheory}), and
for the number of distinct subtree sizes along a root path in heavy-light decomposition ($\Oh{\log n}$ light edges). A
second family of small answers comes from sums: the harmonic sum $\frac n1 + \frac n2 + \dots + \frac nn \approx n \ln n$
means that ``for each $d$, visit every multiple of $d$'' costs only $\Oh{n \log n}$ in total.

\begin{trick}[Swap the loops]
``For each $x$, iterate over its divisors'' costs $\Oh{n\sqrt n}$ if done naively, but ``for each $d$, iterate over its
multiples'' touches the same pairs in $\Oh{n \log n}$ by the harmonic sum. Whenever a double loop's inner range depends on
the outer variable, ask whether the other order is cheaper.
\end{trick}

% ======================================================================================================================
\section[Heavy and light]{Heavy and light: the square-root split \hfill \plat}\label{sec:heavylight}

When some objects are big and others small, a single method is often bad for one of the groups: iterating over an object
costs its size, and precomputing an answer for an object costs a fixed amount. Split them at a threshold $B$: handle
\emph{light} objects (size below $B$) by brute force and \emph{heavy} ones (size at least $B$) by precomputation. If the
total size is $S$, there are at most $S / B$ heavy objects, and choosing $B = \sqrt S$ balances the two costs.

\begin{problem}{Set intersections}
There are $n$ sets of integers with total size $S \le 2\cdot10^5$, and $q$ queries ``how many elements do sets $x$ and $y$
have in common?''.
\solution If one of the two sets is light, iterate over it and look each element up in a hash set of the other: $\Oh{B}$
per query. If both are heavy, answer from a table: for each heavy set $h$, mark its elements and scan every other set once,
recording $|h \cap s|$ for all $s$; that is $\Oh{S}$ per heavy set and $\Oh{S \cdot S / B}$ in total. With $B = \sqrt S$
everything costs $\Oh{(S + q)\sqrt S}$ (Figure~\ref{fig:heavylight}).
\end{problem}

\begin{figure}[htbp]
\centering
\begin{tikzpicture}[x=0.55cm, y=0.3cm]
  \foreach \h [count=\i from 0] in {9.5, 8.2, 7.1, 3.1, 2.6, 2.2, 1.9, 1.6, 1.4, 1.2, 1.0, 0.9, 0.8, 0.7, 0.6, 0.5, 0.5, 0.4} {
    \pgfmathsetmacro{\st}{\h > 4 ? 1 : 0}
    \ifdim\st pt>0.5pt \fill[fgverm!45, draw=fgverm] (\i, 0) rectangle ++(0.8, \h);
    \else \fill[fgblue!30, draw=fgblue] (\i, 0) rectangle ++(0.8, \h); \fi
  }
  \draw[line width=1.2pt, fgink, dashed] (-0.3, 4) -- (18.3, 4) node[right, font=\sffamily\scriptsize] {$B = \sqrt S$};
  \draw[->, line width=0.6pt, fgink!70] (-0.3, 0) -- (18.5, 0) node[right, lbl] {sets, by size};
  \node[font=\sffamily\scriptsize, text=fgverm, align=center] at (1.4, 11.3) {heavy: at most $S/B$ of them\\precompute all answers};
  \node[font=\sffamily\scriptsize, text=fgblue, align=center] at (11.5, 6.6) {light: each smaller than $B$\\brute force per query};
\end{tikzpicture}
\caption{The square-root split. Few objects can be large (their sizes add up to at most $S$), so they can afford expensive
precomputation; the many small ones are cheap to handle directly.}
\label{fig:heavylight}
\end{figure}

The split appears in graphs (vertices of degree above $\sqrt m$ are heavy; there are at most $2\sqrt m$ of them, which gives
the $\Oh{m\sqrt m}$ triangle-counting algorithm), in strings (patterns longer than $\sqrt L$ are few), and in queries
(rebuild a structure every $\sqrt q$ updates and brute force the pending ones, Chapter~\ref{ch:advds}). It rarely gives the
fastest possible algorithm, but it often turns a problem nobody can see how to solve into one that is merely careful work.

% ======================================================================================================================
\section[Differences instead of values]{Differences instead of values \hfill \silver}\label{sec:differences}

Operations on ranges of an array are awkward to reason about directly, but in the array of \emph{differences}
$d_i = b_i - b_{i-1}$ they become local: adding 1 to $b_l, \dots, b_r$ changes only $d_l$ (up by one) and $d_{r+1}$ (down by
one). Many ``minimum number of range operations'' problems collapse once rewritten this way.

\begin{problem}{Range increments}
Starting from an array of zeros, one operation adds 1 to every element of a subarray. What is the minimum number of
operations that produce a given target $b$ ($b_i \ge 0$)?
\solution With $b_0 = 0$ before the array, each operation raises one difference by one and lowers another by one (or none,
at the end). The positive differences sum to $\sum_i \max(0, b_i - b_{i-1})$, and every operation raises that sum by at most
one, so at least that many operations are needed. It is also enough: peel the target into horizontal layers
(Figure~\ref{fig:differences}); each maximal horizontal run in a layer is one operation, and each run starts exactly at a
rise.
\end{problem}

\begin{figure}[htbp]
\centering
\begin{tikzpicture}[x=0.9cm, y=0.55cm]
  % target 2 4 1 3 3 0 decomposed into maximal runs per layer
  \foreach \l/\a/\b/\c in {1/0/5/fgblue, 2/0/2/fggreen, 2/3/5/fggreen, 3/1/2/fgorange, 3/3/5/fgorange, 4/1/2/fgpink} {
    \fill[\c!35, draw=\c!80!black, line width=0.6pt, rounded corners=1pt] (\a+0.04, \l-0.96) rectangle (\b-0.04, \l-0.04);
  }
  \foreach \v [count=\i from 0] in {2,4,1,3,3,0} {
    \node[font=\sffamily\small] at (\i+0.5, -0.65) {\v};
  }
  \node[lbl, anchor=east] at (-0.1, -0.65) {$b$};
  \foreach \i/\d in {0/+2, 1/+2, 3/+2} \node[font=\sffamily\scriptsize\bfseries, text=fgverm] at (\i+0.5, -1.55) {\d};
  \node[lbl, anchor=east] at (-0.1, -1.55) {rises};
  \draw[line width=0.6pt, fgink!70] (0, 0) -- (6, 0);
  \node[notel, anchor=west, text width=5.8cm] at (7, 2) {Each coloured bar is one operation (a maximal run inside one
    layer). There are $1 + 2 + 2 + 1 = 6$ of them, and every bar starts at a position where $b$ rises: the rises
    $2 + 2 + 2 = 6$ give the answer directly.};
\end{tikzpicture}
\caption{Target $2, 4, 1, 3, 3, 0$ cut into horizontal layers. The number of runs equals the sum of the positive
differences $b_i - b_{i-1}$, and no sequence of range increments can do better.}
\label{fig:differences}
\end{figure}

\begin{code}
void solve() {
    int n;
    cin >> n;
    ll prev = 0, ops = 0;
    rep(i, n) {
        ll b;
        cin >> b;
        ops += max(0LL, b - prev);  // a rise of d starts d new ranges here
        prev = b;
    }
    print(ops);
}
\end{code}

Differences also turn ``add $x$ to a range'' into two point updates (the difference array of Chapter~\ref{ch:prefix}), turn
arithmetic progressions into constants (take differences twice), and make ``all elements equal'' into ``all differences
zero''. The inverse operation, prefix sums, turns a range query into two point queries. Keeping both views in mind is
what makes range problems tractable.

% ======================================================================================================================
\section[Reductions to known problems]{Reductions to known problems \hfill \gold}\label{sec:reductions}

The fastest way to solve a problem is to notice that it has already been solved. A large share of contest problems are
classic problems wearing a costume, and part of training is building a mental catalogue of costumes.

\begin{center}
\small
\begin{tabularx}{\linewidth}{Lp{6.2cm}}
\toprule
\thd{The problem says\ldots} & \thd{\ldots and is really} \\
\midrule
fewest non-increasing sequences covering an array & longest strictly increasing subsequence (Dilworth), Ch.~\ref{ch:dp} \\
boxes that nest inside each other, maximise a chain & sort by one side, LIS on the other \\
assign workers to jobs, each to at most one & bipartite matching or min-cost flow, Ch.~\ref{ch:flows} \\
remove fewest vertices so no conflicting pair remains & vertex cover; on bipartite graphs, a matching (K\H{o}nig) \\
variables with ``$x_u - x_v \le c$'' constraints & shortest paths with negative edges, Ch.~\ref{ch:shortest} \\
each variable true or false, clauses of two literals & 2-SAT via strongly connected components, Ch.~\ref{ch:connectivity} \\
minimum cost to make a sequence non-decreasing & slope trick, Ch.~\ref{ch:dpopt} \\
minimise a maximum edge on a path & minimum spanning tree (its path is optimal), Ch.~\ref{ch:dsu} \\
states and moves, fewest moves & BFS on an implicit graph, Ch.~\ref{ch:graphs} \\
choose a subset with a sum constraint, maximise value & knapsack, Ch.~\ref{ch:dp} \\
\bottomrule
\end{tabularx}
\end{center}

\begin{tip}
When a problem looks new, restate it with all names removed, then ask: what is the underlying object (sequence, graph,
set system), what is the constraint, what is optimised? Compare that skeleton with the classics above and with the
recognition table of Chapter~\ref{ch:craft}.
\end{tip}

\begin{summary}[Chapter summary]
\begin{itemize}
  \item Count by \textbf{contribution} and by \textbf{complement}; fix the part that makes the rest independent.
  \item Reorder time: go \textbf{offline}, run it \textbf{backwards}, or \textbf{sweep} over sorted events.
  \item Prove impossibility with \textbf{invariants} and termination with \textbf{monovariants}; prove existence with
        \textbf{pigeonhole}, \textbf{extremal} objects and \textbf{reflection}.
  \item Exploit \textbf{small answers} (divisor chains, harmonic sums), the \textbf{heavy/light} split, and
        \textbf{differences} for range operations; above all, look for a known problem in disguise.
\end{itemize}
\end{summary}

\begin{practice}
\ppgroup{Contribution and complement}
\begin{pplist}
  \item Sum of Two Values
  \item Subarray Sums II
  \item Nearest Smaller Values
  \item Distinct Numbers
  \item Maximum Subarray Sum
\end{pplist}
\ppgroup{Offline, reversed, sweeps}
\begin{pplist}
  \item Restaurant Customers
  \item Room Allocation
  \item Network Breakdown
  \item Nested Ranges Check
  \item Distinct Values Queries
\end{pplist}
\ppgroup{Invariants, pigeonhole, differences}
\begin{pplist}
  \item Subarray Divisibility
  \item Two Sets
  \item Coin Piles
  \item Increasing Array
  \item Collecting Numbers
  \item Bracket Sequences I
\end{pplist}
\end{practice}
